Optimal stopping compares the payoff available now with the value of continuing and then acting optimally. The state must include the remaining opportunities, any offer that can be recalled, and the posterior for an unknown parameter. Do not charge costs already paid a second time. Compare feasible actions at each observation before averaging over observations; equality may allow several optimal decisions.

Backward induction verifies finite-horizon policies. With unlimited opportunities, establish a finite value, justify the limit of the recursion, and check that the proposed policy attains it. Discounting can provide a uniform bound on truncation error. Replication and superhedging also require nodewise checks of self-financing, borrowing and short-selling restrictions, not just comparison of terminal payoffs.
\par\smallskip
